\par\Needspace{6\baselineskip}\medskip
\hypertarget{native-208-article154}{}\pdfbookmark[2]{154.}{article-154}
\begin{center}154.\end{center}
Der Vollständigkeit halber stellen wir noch eine \textit{dritte Hypothese} auf. Dabei übernehmen wir wieder die in der zweiten Hypothese ermittelten Werte von $P',Q'$ als Werte von $P,Q$. Setzt man also
\[
\begin{aligned}
x&=\log P=0{,}079\,0167\\
y&=\log Q=8{,}547\,6110,
\end{aligned}
\]
\Needspace{14\baselineskip}
so erhält man die folgenden wichtigsten Ergebnisse der Rechnung:
\begin{center}\small\setlength{\tabcolsep}{6pt}
\begin{tabular}{@{}r|r@{}}
$\omega=13^{0}\,15'\,38''{,}39$ & $\zeta''=210^{0}\,8'\,25''{,}65$\\[.5em]
$\omega+\sigma=13\quad38\quad51{,}51$ & $\log r=0{,}330\,7640$\\[.5em]
$\log Qc\sin\omega=0{,}598\,9542$ & $\log r''=0{,}322\,2239$\\[.5em]
$z=14^{0}\,33'\,19''{,}50$ & $\tfrac12(u''+u)=205^{0}\,22'\,14''{,}57$\\[.5em]
$\log r'=0{,}325\,9878$ & $\tfrac12(u''-u)=-3\quad14\quad4{,}78$\\[.5em]
$\log\dfrac{n'r'}{n}=0{,}667\,5154$ & $2f'=7\quad34\quad53{,}73$\\[.7em]
$\log\dfrac{n'r'}{n''}=0{,}588\,4987$ & $2f=3\quad29\quad0{,}39$\\[.7em]
$\zeta=203^{0}\,16'\,38''{,}41$ & $2f''=4\quad5\quad53{,}34$
\end{tabular}
\end{center}

Alle diese Zahlen unterscheiden sich so wenig von denen der zweiten Hypothese, dass wir mit Sicherheit schließen können, dass die dritte Hypothese keiner weiteren Korrektur bedarf\footnote{Führte man die Rechnung wie in den vorangegangenen Hypothesen zu Ende, so erhielte man $X=0$, $Y=-0{,}000\,0003$. Der letztere Wert kann als verschwindend angesehen werden und geht kaum über die Unsicherheit hinaus, die der letzten Dezimalstelle stets anhaftet.}. Wir können daher zur Bestimmung der Elemente selbst aus $2f',r,r'',\theta'$ übergehen. Auf ihre Wiedergabe verzichten wir hier, da sie bereits oben in Artikel~97 als Beispiel ausführlich dargestellt wurde. Es bleibt nur noch, die Lage der Bahnebene nach Artikel~149 zu berechnen und die Epoche auf den Beginn des Jahres 1805 zu verlegen. Dieser Rechnung sind die folgenden Zahlen zugrunde zu legen:
\SourceLeaf{209}
\[
\begin{aligned}
AD'-\zeta&=9^{0}\,55'\,51''{,}41\\
\tfrac12(\gamma+u)&=202\quad18\quad13{,}855\\
\tfrac12(\gamma-u)&=-6\quad18\quad5{,}495,
\end{aligned}
\]
daraus leiten wir ab
\[
\begin{aligned}
\tfrac12(g+h)&=196^{0}\,43'\,14''{,}62\\
\tfrac12(g-h)&=-4\quad37\quad24{,}41\\
\tfrac12i&=6\quad33\quad22{,}05.
\end{aligned}
\]
Damit ergibt sich $h=201^{0}\,20'\,39''{,}03$, also $\AscendingNode=l-h=171^{0}\,7'\,48''{,}73$. Ferner ergibt sich $g=192^{0}\,5'\,50''{,}21$. Da die wahre Anomalie für die erste Beobachtung in Artikel~97 zu $=310^{0}\,55'\,29''{,}64$ gefunden wurde, ist der Abstand des Perihels vom aufsteigenden Knoten in der Bahn $=241^{0}\,10'\,20''{,}57$ und die Länge des Perihels $=52^{0}\,18'\,9''{,}30$. Schließlich beträgt die Bahnneigung $=13^{0}\,6'\,44''{,}10$. --- Will man bei derselben Rechnung lieber vom dritten Ort ausgehen, so hat man
\[
\begin{aligned}
A''D'-\zeta''&=24^{0}\,18'\,35''{,}25\\
\tfrac12(\gamma''+u'')&=196\quad24\quad54{,}98\\
\tfrac12(\gamma''-u'')&=-5\quad43\quad14{,}81.
\end{aligned}
\]
Daraus folgt
\[
\begin{aligned}
\tfrac12(g''+h'')&=211^{0}\,24'\,32''{,}44\\
\tfrac12(g''-h'')&=-11\quad43\quad48{,}49\\
\tfrac12i&=6\quad33\quad22{,}05,
\end{aligned}
\]
und daraus die Länge des aufsteigenden Knotens $=l''-h''=171^{0}\,7'\,48''{,}72$, die Länge des Perihels $=52^{0}\,18'\,9''{,}30$ und die Bahnneigung $=13^{0}\,6'\,44''{,}10$, ganz wie zuvor.

Das Zeitintervall von der letzten Beobachtung bis zum Beginn des Jahres 1805 beträgt $64{,}614\,102$ Tage. Ihm entspricht die mittlere heliozentrische Bewegung $53293''{,}65=14^{0}\,48'\,13''{,}65$. Damit ergibt sich die Epoche der mittleren Anomalie für den Beginn des Jahres 1805 am Pariser Meridian $=349^{0}\,34'\,12''{,}38$ und die Epoche der mittleren Länge $=41^{0}\,52'\,21''{,}68$.
